\section{Seguridad e infraestructura}

\subsection{Registro y observabilidad}

Antes de desplegar una reward learned policy, es necesario garantizar una cadena completa de registro y observabilidad. Las dimensiones clave incluyen:

\begin{itemize}
    \item \textbf{Reward Trace}: registra la puntuación del reward model, la predicción de la baseline policy, el KL penalty
    \item \textbf{Action Trace}: registra la acción elegida por la policy, la distribución de probabilidad, la temperature
    \item \textbf{Safety Signals}: detecta contenido inapropiado, hallucination, comportamiento repetitivo
\end{itemize}

\begin{warningbox}{Sin registro, los resultados no se pueden corregir}
En los sistemas reales, lo más peligroso del reward hacking no es que la salida del modelo se desvíe, sino que el equipo simplemente no puede ver dónde se «quiebra». La práctica habitual es escribir el reward model score en la trace table y hacer un roll back automático cuando no se alcanza el umbral.
\end{warningbox}

\subsection{De la simulación a la realidad}

El entrenamiento de reward learning suele realizarse en un entorno de simulación, por lo que antes de desplegarlo en el mundo real es necesario hacer un «complemento de sim-to-real»:

\begin{itemize}
    \item \textbf{Domain randomization}: agregar iluminación, ruido y perturbaciones físicas durante el entrenamiento
    \item \textbf{Reward smoothing}: agregar dropout a la reward signal o mezclarla con etiquetas humanas reales, para evitar el overfit
    \item \textbf{Distillation+Calibration}: destilar la large policy en un modelo pequeño y, al mismo tiempo, calibrar el reward model para contrarrestar el distribution shift
\end{itemize}

\begin{knowledgebox}{Caso de sim-to-real}
En una tarea de grasping con robots, la reward learned policy supera el 95\% en simulación, pero al desplegarla directamente en una pinza real cae hasta el 45\%. La solución es mezclar la reward signal recolectada en la simulación con una pequeña cantidad de human data real, y agregar perturbaciones de deformación para hacer fine-tuning.
\end{knowledgebox}

\subsection{Resumen de la sección}

El despliegue seguro exige un sistema completo de registro, revisión automática y calibración de sim-to-real. Cada paso del reward model debe poder rastrearse y revertirse.
